\subsection{Produktzeichen}
\label{subsec:fundamentals_product_notation}

Analog zum Summenzeichen fasst das Produktzeichen zu multiplizierende
Folgenglieder kompakt zusammen.

\begin{DefinitionBox}{Produktschreibweise}{fundamentals_product_notation}
  Die \emph{Produktschreibweise} \EnglishTerm{product notation}
  \[
    \prod_{k=m}^{n}a_k
      =a_m\cdot a_{m+1}\cdots a_n
  \]
  fasst die Glieder $a_m$ bis $a_n$ zu einem Produkt zusammen. Wie beim
  Summenzeichen ist $k$ der Laufindex; $m$ und $n$ sind die untere und obere
  Grenze.
\end{DefinitionBox}

\begin{ExampleBox}{Produkte auswerten}{fundamentals_evaluating_products}
  \begin{enumerate}[label=(\alph*),leftmargin=2em,itemsep=0.8em]
    \item\leavevmode
      \[
        \prod_{k=1}^{5}k
          =\Cloze{12cm}{1\cdot2\cdot3\cdot4\cdot5=120}
      \]
    \item\leavevmode
      \[
        \prod_{k=1}^{4}k^2
          =\Cloze{12cm}{1^2\cdot2^2\cdot3^2\cdot4^2
            =1\cdot4\cdot9\cdot16=576}
      \]
    \item\leavevmode
      \[
        \prod_{k=1}^{3}\frac1k
          =\Cloze{12cm}{1\cdot\frac12\cdot\frac13=\frac16}
      \]
  \end{enumerate}
\end{ExampleBox}
